\let\textcircled=\pgftextcircled
\chapter{Pr\'esentation de la structure d'accueil}
\label{chap:structure}

\textit{\initial{C}e chapitre pr\'esente l'organisme au sein duquel le stage s'est
d\'eroul\'e, la mission qui y a \'et\'e confi\'ee et le p\'erim\`etre op\'erationnel dans lequel
le syst\`eme con\c cu doit s'ins\'erer.}

\mtcselectlanguage{french}
\minitoc
\newpage

\section{L'Agence Nationale des Technologies de l'Information et de la Communication}
% ---------------------------------------------------------------------------
% A COMPLETER. Elements a reunir aupres de la structure :
%   — texte de creation (decret) et missions statutaires ;
%   — tutelle ministerielle ;
%   — effectif, implantation, budget si communicable.
% ---------------------------------------------------------------------------
[Historique, texte de cr\'eation, missions statutaires de l'ANTIC.]

\subsection{Missions}
[\'Enum\'erer les missions : r\'egulation, certification \'electronique, veille,
r\'eponse aux incidents, sensibilisation.]

\subsection{Organisation}
[Organigramme et positionnement de la direction dont d\'epend le CIRT.]

\begin{figure}[H]
    \centering
    % \includegraphics[width=0.85\linewidth]{fig01/organigramme.png}
    \caption{Organigramme de l'ANTIC}
    \label{fig:organigramme}
\end{figure}

\section{Le CIRT national}
\subsection{Mandat et p\'erim\`etre}
[Champ de comp\'etence : administrations, op\'erateurs d'importance vitale, secteur
bancaire, signalements citoyens.]

\subsection{Organisation de la r\'eponse aux incidents}
[Cha\^ine actuelle : r\'eception du signalement, qualification, escalade,
coordination, cl\^oture. Pr\'eciser les niveaux d'astreinte.]

\subsection{Outillage existant}
[Sources de d\'etection en place : SIEM, sondes r\'eseau, EDR, remont\'ees
partenaires. Cet inventaire conditionne les adaptateurs \`a d\'evelopper --- voir
la figure~\ref{fig:composants}.]

\section{Cadre et d\'eroulement du stage}
\subsection{Objet de la mission}
[Intitul\'e exact de la mission tel qu'il figure sur la convention de stage.]

\subsection{Chronologie des travaux}
% Remplacer les intitules et les bornes par le deroule reel.
\begin{table}[H]
\centering
\caption{Chronologie des travaux men\'es durant le stage}
\label{tab:chronologie}
\begin{tabular}{@{}llp{7.5cm}@{}}
\toprule
\textbf{P\'eriode} & \textbf{Phase} & \textbf{Livrable} \\
\midrule
Semaines 1--2 & Immersion    & Prise de connaissance des proc\'edures du CIRT \\
Semaines 3--4 & Analyse      & Cahier des charges fonctionnel \\
Semaines 5--7 & Conception   & Mod\`eles UML et architecture \\
Semaines 8--13 & R\'ealisation & Plateforme et suite de tests \\
Semaine 14    & Validation   & Recette et r\'edaction du rapport \\
\bottomrule
\end{tabular}
\end{table}

\subsection{Contraintes impos\'ees}
Trois contraintes ont structur\'e l'ensemble du travail et sont reprises comme
exigences non fonctionnelles au chapitre~\ref{chap:analyse} :

\begin{enumerate}
  \item \textbf{Souverainet\'e} --- aucune donn\'ee d'incident ne doit sortir du
        territoire ; le syst\`eme doit fonctionner sans acc\`es Internet ;
  \item \textbf{Opposabilit\'e} --- toute action men\'ee par le syst\`eme doit pouvoir
        \^etre reconstitu\'ee et prouv\'ee a posteriori ;
  \item \textbf{Sobri\'et\'e d'infrastructure} --- le d\'eploiement doit tenir sur un
        serveur unique, sans d\'ependance \`a un service manag\'e externe.
\end{enumerate}

\section{Conclusion du chapitre}
[Synth\`ese : ce que la structure attend du syst\`eme et ce que le chapitre suivant
va confronter \`a l'existant.]
